\rSec1[function.objects]{Function objects}

\rSec2[function.objects.general]{General}

\pnum
A \defnx{function object type}{function object!type} is an object
type\iref{term.object.type} that can be the type of the
\grammarterm{postfix-expression}
in a function call\iref{expr.call,over.match.call}.
\begin{footnote}
Such a type is a function
pointer or a class type which has a member \tcode{operator()} or a class type
which has a conversion to a pointer to function.
\end{footnote}
A \defn{function object} is an
object of a function object type. In the places where one would expect to pass a
pointer to a function to an algorithmic template\iref{algorithms}, the
interface is specified to accept a function object. This not only makes
algorithmic templates work with pointers to functions, but also enables them to
work with arbitrary function objects.

\rSec2[functional.syn]{Header \tcode{<functional>} synopsis}

\indexheader{functional}%
\indexlibraryglobal{unwrap_ref_decay}%
\indexlibraryglobal{unwrap_ref_decay_t}%
\begin{codeblock}
namespace std {
  // \ref{func.invoke}, invoke
  template<class F, class... Args>
    constexpr invoke_result_t<F, Args...> invoke(F&& f, Args&&... args)             // freestanding
      noexcept(is_nothrow_invocable_v<F, Args...>);

  template<class R, class F, class... Args>
    constexpr R invoke_r(F&& f, Args&&... args)                                     // freestanding
      noexcept(is_nothrow_invocable_r_v<R, F, Args...>);

  // \ref{refwrap}, \tcode{reference_wrapper}
  template<class T> class reference_wrapper;                                        // freestanding

  template<class T> constexpr reference_wrapper<T> ref(T&) noexcept;                // freestanding
  template<class T> constexpr reference_wrapper<const T> cref(const T&) noexcept;   // freestanding
  template<class T> void ref(const T&&) = delete;                                   // freestanding
  template<class T> void cref(const T&&) = delete;                                  // freestanding

  template<class T>
    constexpr reference_wrapper<T> ref(reference_wrapper<T>) noexcept;              // freestanding
  template<class T>
    constexpr reference_wrapper<const T> cref(reference_wrapper<T>) noexcept;       // freestanding

  // \ref{refwrap.common.ref}, \tcode{common_reference} related specializations
  template<class R, class T, template<class> class RQual, template<class> class TQual>
    requires @\seebelow@
  struct basic_common_reference<R, T, RQual, TQual>;

  template<class T, class R, template<class> class TQual, template<class> class RQual>
    requires @\seebelow@
  struct basic_common_reference<T, R, TQual, RQual>;

  // \ref{arithmetic.operations}, arithmetic operations
  template<class T = void> struct plus;                                             // freestanding
  template<class T = void> struct minus;                                            // freestanding
  template<class T = void> struct multiplies;                                       // freestanding
  template<class T = void> struct divides;                                          // freestanding
  template<class T = void> struct modulus;                                          // freestanding
  template<class T = void> struct negate;                                           // freestanding
  template<> struct plus<void>;                                                     // freestanding
  template<> struct minus<void>;                                                    // freestanding
  template<> struct multiplies<void>;                                               // freestanding
  template<> struct divides<void>;                                                  // freestanding
  template<> struct modulus<void>;                                                  // freestanding
  template<> struct negate<void>;                                                   // freestanding

  // \ref{comparisons}, comparisons
  template<class T = void> struct equal_to;                                         // freestanding
  template<class T = void> struct not_equal_to;                                     // freestanding
  template<class T = void> struct greater;                                          // freestanding
  template<class T = void> struct less;                                             // freestanding
  template<class T = void> struct greater_equal;                                    // freestanding
  template<class T = void> struct less_equal;                                       // freestanding
  template<> struct equal_to<void>;                                                 // freestanding
  template<> struct not_equal_to<void>;                                             // freestanding
  template<> struct greater<void>;                                                  // freestanding
  template<> struct less<void>;                                                     // freestanding
  template<> struct greater_equal<void>;                                            // freestanding
  template<> struct less_equal<void>;                                               // freestanding

  // \ref{comparisons.three.way}, class \tcode{compare_three_way}
  struct compare_three_way;                                                         // freestanding

  // \ref{logical.operations}, logical operations
  template<class T = void> struct logical_and;                                      // freestanding
  template<class T = void> struct logical_or;                                       // freestanding
  template<class T = void> struct logical_not;                                      // freestanding
  template<> struct logical_and<void>;                                              // freestanding
  template<> struct logical_or<void>;                                               // freestanding
  template<> struct logical_not<void>;                                              // freestanding

  // \ref{bitwise.operations}, bitwise operations
  template<class T = void> struct bit_and;                                          // freestanding
  template<class T = void> struct bit_or;                                           // freestanding
  template<class T = void> struct bit_xor;                                          // freestanding
  template<class T = void> struct bit_not;                                          // freestanding
  template<> struct bit_and<void>;                                                  // freestanding
  template<> struct bit_or<void>;                                                   // freestanding
  template<> struct bit_xor<void>;                                                  // freestanding
  template<> struct bit_not<void>;                                                  // freestanding

  // \ref{func.identity}, identity
  struct identity;                                                                  // freestanding

  // \ref{func.not.fn}, function template \tcode{not_fn}
  template<class F> constexpr @\unspec@ not_fn(F&& f);                            // freestanding
  template<auto f> constexpr @\unspec@ not_fn() noexcept;                         // freestanding

  // \ref{func.bind.partial}, function templates \tcode{bind_front} and \tcode{bind_back}
  template<class F, class... Args>
    constexpr @\unspec@ bind_front(F&&, Args&&...);                               // freestanding
  template<auto f, class... Args>
    constexpr @\unspec@ bind_front(Args&&...);                                    // freestanding
  template<class F, class... Args>
    constexpr @\unspec@ bind_back(F&&, Args&&...);                                // freestanding
  template<auto f, class... Args>
    constexpr @\unspec@ bind_back(Args&&...);                                     // freestanding

  // \ref{func.bind}, bind
  template<class T> struct is_bind_expression;                                      // freestanding
  template<class T>
    constexpr bool @\libglobal{is_bind_expression_v}@ =                                           // freestanding
      is_bind_expression<T>::value;
  template<class T> struct is_placeholder;                                          // freestanding
  template<class T>
    constexpr int @\libglobal{is_placeholder_v}@ =                                                // freestanding
      is_placeholder<T>::value;

  template<class F, class... BoundArgs>
    constexpr @\unspec@ bind(F&&, BoundArgs&&...);                                // freestanding
  template<class R, class F, class... BoundArgs>
    constexpr @\unspec@ bind(F&&, BoundArgs&&...);                                // freestanding

  namespace placeholders {
    // \tcode{\placeholder{M}} is the \impldef{number of placeholders for bind expressions} number of placeholders
    @\seebelownc@ _1;                                                                   // freestanding
    @\seebelownc@ _2;                                                                   // freestanding
               .
               .
               .
    @\seebelownc@ _@\placeholdernc{M}@;                                                                   // freestanding
  }

  // \ref{func.memfn}, member function adaptors
  template<class R, class T>
    constexpr @\unspec@ mem_fn(R T::*) noexcept;                                  // freestanding

  // \ref{func.wrap}, polymorphic function wrappers
  // \ref{func.wrap.badcall}, class \tcode{bad_function_call}
  class bad_function_call;

  // \ref{func.wrap.func}, class template \tcode{function}
  template<class> class function;       // \notdef
  template<class R, class... ArgTypes> class function<R(ArgTypes...)>;

  // \ref{func.wrap.func.alg}, \tcode{function} specialized algorithms
  template<class R, class... ArgTypes>
    void swap(function<R(ArgTypes...)>&, function<R(ArgTypes...)>&) noexcept;

  // \ref{func.wrap.func.nullptr}, \tcode{function} null pointer comparison operator functions
  template<class R, class... ArgTypes>
    bool operator==(const function<R(ArgTypes...)>&, nullptr_t) noexcept;

  // \ref{func.wrap.move}, move-only wrapper
  template<class... S> class move_only_function;                        // \notdef
  template<class R, class... ArgTypes>
    class move_only_function<R(ArgTypes...) @\cv{}\itcorr[-1]@ @\placeholdernc{ref}@ noexcept(@\placeholdernc{noex}@)>;     // \seebelow

  // \ref{func.wrap.copy}, copyable wrapper
  template<class... S> class copyable_function;                         // \notdef
  template<class R, class... ArgTypes>
    class copyable_function<R(ArgTypes...) @\cv{}\itcorr[-1]@ @\placeholdernc{ref}@ noexcept(@\placeholdernc{noex}@)>;      // \seebelow

  // \ref{func.wrap.ref}, non-owning wrapper
  template<class... S> class function_ref;                              // freestanding, \notdef
  template<class R, class... ArgTypes>
    class function_ref<R(ArgTypes...) @\cv{}\itcorr[-1]@ noexcept(@\placeholdernc{noex}@)>;               // freestanding, \seebelow

  // \ref{func.search}, searchers
  template<class ForwardIterator1, class BinaryPredicate = equal_to<>>
    class default_searcher;                                                         // freestanding

  template<class RandomAccessIterator,
           class Hash = hash<typename iterator_traits<RandomAccessIterator>::value_type>,
           class BinaryPredicate = equal_to<>>
    class boyer_moore_searcher;

  template<class RandomAccessIterator,
           class Hash = hash<typename iterator_traits<RandomAccessIterator>::value_type>,
           class BinaryPredicate = equal_to<>>
    class boyer_moore_horspool_searcher;

  // \ref{unord.hash}, class template \tcode{hash}
  template<class T>
    struct hash;                                                                    // freestanding

  namespace ranges {
    // \ref{range.cmp}, concept-constrained comparisons
    struct equal_to;                                                                // freestanding
    struct not_equal_to;                                                            // freestanding
    struct greater;                                                                 // freestanding
    struct less;                                                                    // freestanding
    struct greater_equal;                                                           // freestanding
    struct less_equal;                                                              // freestanding
  }

  template<class Fn, class... Args>
    concept @\defexposconceptnc{callable}@ =                                                  // \expos
      requires (Fn&& fn, Args&&... args) {
        std::forward<Fn>(fn)(std::forward<Args>(args)...);
      };

  template<class Fn, class... Args>
    concept @\defexposconceptnc{nothrow-callable}@ =                                          // \expos
      @\exposconcept{callable}@<Fn, Args...> &&
      requires (Fn&& fn, Args&&... args) {
        { std::forward<Fn>(fn)(std::forward<Args>(args)...) } noexcept;
      };

  template<class Fn, class... Args>
    using @\exposidnc{call-result-t}@ = decltype(declval<Fn>()(declval<Args>()...));  // \expos

  template<const auto& T>
    using @\exposidnc{decayed-typeof}@ = decltype(auto(T));                           // \expos
}
\end{codeblock}

\pnum
\begin{example}
If a \Cpp{} program wants to have a by-element addition of two vectors \tcode{a}
and \tcode{b} containing \tcode{double} and put the result into \tcode{a},
it can do:

\begin{codeblock}
transform(a.begin(), a.end(), b.begin(), a.begin(), plus<double>());
\end{codeblock}
\end{example}

\pnum
\begin{example}
To negate every element of \tcode{a}:

\begin{codeblock}
transform(a.begin(), a.end(), a.begin(), negate<double>());
\end{codeblock}

\end{example}

\rSec2[func.def]{Definitions}

\rSec2[func.require]{Requirements}

\rSec2[func.invoke]{\tcode{invoke} functions}

\rSec2[refwrap]{Class template \tcode{reference_wrapper}}

\rSec2[arithmetic.operations]{Arithmetic operations}

\rSec2[comparisons]{Comparisons}

\rSec2[range.cmp]{Concept-constrained comparisons}

\rSec2[logical.operations]{Logical operations}

\rSec2[bitwise.operations]{Bitwise operations}

\rSec2[func.identity]{Class \tcode{identity}}

\rSec2[func.not.fn]{Function template \tcode{not_fn}}

\rSec2[func.bind.partial]{Function templates \tcode{bind_front} and \tcode{bind_back}}

\rSec2[func.bind]{Function object binders}%

\rSec2[func.memfn]{Function template \tcode{mem_fn}}%

\rSec2[func.wrap]{Polymorphic function wrappers}%

\rSec3[func.wrap.general]{General}%
\indextext{function object!wrapper|(}

\pnum
Subclause \ref{func.wrap} describes polymorphic wrapper classes that
encapsulate arbitrary callable objects.

\pnum
Let \tcode{t} be an object of a type that is a specialization of
\tcode{function}, \tcode{copyable_function}, or \tcode{move_only_function},
such that the target object \tcode{x} of \tcode{t} has a type that
is a specialization of
\tcode{function}, \tcode{copyable_function}, or \tcode{move_only_function}.
Each argument of the
invocation of \tcode{x} evaluated as part of the invocation of \tcode{t}
may alias an argument in the same position in the invocation of \tcode{t} that
has the same type, even if the corresponding parameter is not of reference type.
\begin{example}
\begin{codeblock}
move_only_function<void(T)>
  f{copyable_function<void(T)>{[](T) {}}};
T t;
f(t);                               // it is unspecified how many copies of \tcode{T} are made
\end{codeblock}
\end{example}

\pnum
\recommended
Implementations should avoid double wrapping when
constructing polymorphic wrappers from one another.

\rSec3[func.wrap.badcall]{Class \tcode{bad_function_call}}%
\indexlibraryglobal{bad_function_call}%

\pnum
An exception of type \tcode{bad_function_call} is thrown by
\tcode{function::operator()}\iref{func.wrap.func.inv}
when the function wrapper object has no target.

\begin{codeblock}
namespace std {
  class bad_function_call : public exception {
  public:
    // see \ref{exception} for the specification of the special member functions
    const char* what() const noexcept override;
  };
}
\end{codeblock}

\indexlibrarymember{what}{bad_function_call}%
\begin{itemdecl}
const char* what() const noexcept override;
\end{itemdecl}

\begin{itemdescr}
\pnum
\returns
An
\impldef{return value of \tcode{bad_function_call::what}} \ntbs{}.
\end{itemdescr}

\rSec3[func.wrap.func]{Class template \tcode{function}}

\rSec4[func.wrap.func.general]{General}
\indexlibraryglobal{function}%

\indexlibrarymember{result_type}{function}%
\begin{codeblock}
namespace std {
  template<class R, class... ArgTypes>
  class function<R(ArgTypes...)> {
  public:
    using result_type = R;

    // \ref{func.wrap.func.con}, construct/copy/destroy
    function() noexcept;
    function(nullptr_t) noexcept;
    function(const function&);
    function(function&&) noexcept;
    template<class F> function(F&&);

    function& operator=(const function&);
    function& operator=(function&&);
    function& operator=(nullptr_t) noexcept;
    template<class F> function& operator=(F&&);
    template<class F> function& operator=(reference_wrapper<F>) noexcept;

    ~function();

    // \ref{func.wrap.func.mod}, function modifiers
    void swap(function&) noexcept;

    // \ref{func.wrap.func.cap}, function capacity
    explicit operator bool() const noexcept;

    // \ref{func.wrap.func.inv}, function invocation
    R operator()(ArgTypes...) const;

    // \ref{func.wrap.func.targ}, function target access
    const type_info& target_type() const noexcept;
    template<class T>       T* target() noexcept;
    template<class T> const T* target() const noexcept;
  };

  template<class R, class... ArgTypes>
    function(R(*)(ArgTypes...)) -> function<R(ArgTypes...)>;

  template<class F> function(F) -> function<@\seebelow@>;
}
\end{codeblock}

\pnum
The \tcode{function} class template provides polymorphic wrappers that
generalize the notion of a function pointer. Wrappers can store, copy,
and call arbitrary callable objects\iref{func.def}, given a call
signature\iref{func.def}.

\pnum
\indextext{callable type}%
A callable type\iref{func.def} \tcode{F}
is \defn{Lvalue-Callable} for argument
types \tcode{ArgTypes}
and return type \tcode{R}
if the expression
\tcode{\placeholdernc{INVOKE}<R>(declval<F\&>(), declval<ArgTypes>()...)},
considered as an unevaluated operand\iref{term.unevaluated.operand}, is
well-formed\iref{func.require}.

\pnum
The \tcode{function} class template is a call
wrapper\iref{func.def} whose call signature\iref{func.def}
is \tcode{R(ArgTypes...)}.

\pnum
\begin{note}
The types deduced by the deduction guides for \tcode{function}
might change in future revisions of \Cpp{}.
\end{note}

\rSec4[func.wrap.func.con]{Constructors and destructor}

\indexlibraryctor{function}%
\begin{itemdecl}
function() noexcept;
\end{itemdecl}

\begin{itemdescr}
\pnum
\ensures
\tcode{!*this}.
\end{itemdescr}

\indexlibraryctor{function}%
\begin{itemdecl}
function(nullptr_t) noexcept;
\end{itemdecl}

\begin{itemdescr}
\pnum
\ensures
\tcode{!*this}.
\end{itemdescr}

\indexlibraryctor{function}%
\begin{itemdecl}
function(const function& f);
\end{itemdecl}

\begin{itemdescr}
\pnum
\ensures
\tcode{!*this} if \tcode{!f}; otherwise,
the target object of \tcode{*this} is a copy of the target object of \tcode{f}.

\pnum
\throws
Nothing if \tcode{f}'s target is
a specialization of \tcode{reference_wrapper} or
a function pointer. Otherwise, may throw \tcode{bad_alloc}
or any exception thrown by the copy constructor of the stored callable object.

\pnum
\recommended
Implementations should avoid the use of
dynamically allocated memory for small callable objects, for example, where
\tcode{f}'s target is an object holding only a pointer or reference
to an object and a member function pointer.
\end{itemdescr}

\indexlibraryctor{function}%
\begin{itemdecl}
function(function&& f) noexcept;
\end{itemdecl}

\begin{itemdescr}
\pnum
\ensures
If \tcode{!f}, \tcode{*this} has no target;
otherwise, the target of \tcode{*this} is equivalent to
the target of \tcode{f} before the construction, and
\tcode{f} is in a valid state with an unspecified value.

\pnum
\recommended
Implementations should avoid the use of
dynamically allocated memory for small callable objects, for example,
where \tcode{f}'s target is an object holding only a pointer or reference
to an object and a member function pointer.
\end{itemdescr}

\indexlibraryctor{function}%
\begin{itemdecl}
template<class F> function(F&& f);
\end{itemdecl}

\begin{itemdescr}
\pnum
Let \tcode{FD} be \tcode{decay_t<F>}.

\pnum
\constraints
\begin{itemize}
\item
\tcode{is_same_v<remove_cvref_t<F>, function>} is \tcode{false}, and
\item
\tcode{FD} is Lvalue-Callable\iref{func.wrap.func} for argument types
\tcode{ArgTypes...} and return type \tcode{R}.
\end{itemize}

\pnum
\mandates
\begin{itemize}
\item
\tcode{is_copy_constructible_v<FD>} is \tcode{true}, and
\item
\tcode{is_constructible_v<FD, F>} is \tcode{true}.
\end{itemize}

\pnum
\expects
\tcode{FD} meets the \oldconcept{CopyConstructible} requirements.

\pnum
\ensures
\tcode{!*this} is \tcode{true} if any of the following hold:
\begin{itemize}
\item \tcode{f} is a null function pointer value.
\item \tcode{f} is a null member pointer value.
\item \tcode{remove_cvref_t<F>} is
a specialization of the \tcode{function} class template, and
\tcode{!f} is \tcode{true}.
\end{itemize}

\pnum
Otherwise, \tcode{*this} has a target object of type \tcode{FD}
direct-non-list-initialized with \tcode{std::forward<F>(f)}.

\pnum
\throws
Nothing if \tcode{FD} is
a specialization of \tcode{reference_wrapper} or
a function pointer type.
Otherwise, may throw \tcode{bad_alloc} or
any exception thrown by the initialization of the target object.

\pnum
\recommended
Implementations should avoid the use of
dynamically allocated memory for small callable objects, for example,
where \tcode{f} refers to an object holding only a pointer or
reference to an object and a member function pointer.
\end{itemdescr}


\begin{itemdecl}
template<class F> function(F) -> function<@\seebelow@>;
\end{itemdecl}

\begin{itemdescr}
\pnum
\constraints
\tcode{\&F::operator()} is well-formed when treated as an unevaluated operand and either
\begin{itemize}
\item
\tcode{F::operator()} is a non-static member function and
\tcode{decltype(\brk{}\&F::operator())} is either of the form
\tcode{R(G::*)(A...)}~\cv{}~\tcode{\opt{\&}~\opt{noexcept}}
or of the form
\tcode{R(*)(G, A...)~\opt{noexcept}}
for a type \tcode{G}, or
\item
\tcode{F::operator()} is a static member function and
\tcode{decltype(\&F::operator())} is of the form
\tcode{R(*)(A...) \opt{noexcept}}.
\end{itemize}

\pnum
\remarks
The deduced type is \tcode{function<R(A...)>}.

\pnum
\begin{example}
\begin{codeblock}
void f() {
  int i{5};
  function g = [&](double) { return i; };       // deduces \tcode{function<int(double)>}
}
\end{codeblock}
\end{example}
\end{itemdescr}

\indexlibrarymember{operator=}{function}%
\begin{itemdecl}
function& operator=(const function& f);
\end{itemdecl}

\begin{itemdescr}
\pnum
\effects
As if by \tcode{function(f).swap(*this);}

\pnum
\returns
\tcode{*this}.
\end{itemdescr}

\indexlibrarymember{operator=}{function}%
\begin{itemdecl}
function& operator=(function&& f);
\end{itemdecl}

\begin{itemdescr}
\pnum
\effects
Replaces the target of \tcode{*this}
with the target of \tcode{f}.

\pnum
\returns
\tcode{*this}.
\end{itemdescr}

\indexlibrarymember{operator=}{function}%
\begin{itemdecl}
function& operator=(nullptr_t) noexcept;
\end{itemdecl}

\begin{itemdescr}
\pnum
\effects
If \tcode{*this != nullptr}, destroys the target of \keyword{this}.

\pnum
\ensures
\tcode{!(*this)}.

\pnum
\returns
\tcode{*this}.
\end{itemdescr}

\indexlibrarymember{operator=}{function}%
\begin{itemdecl}
template<class F> function& operator=(F&& f);
\end{itemdecl}

\begin{itemdescr}
\pnum
\constraints
\tcode{decay_t<F>} is Lvalue-Callable\iref{func.wrap.func}
for argument types \tcode{ArgTypes...} and return type \tcode{R}.

\pnum
\effects
As if by: \tcode{function(std::forward<F>(f)).swap(*this);}

\pnum
\returns
\tcode{*this}.
\end{itemdescr}

\indexlibrarymember{operator=}{function}%
\begin{itemdecl}
template<class F> function& operator=(reference_wrapper<F> f) noexcept;
\end{itemdecl}

\begin{itemdescr}
\pnum
\effects
As if by: \tcode{function(f).swap(*this);}

\pnum
\returns
\tcode{*this}.
\end{itemdescr}

\indexlibrarydtor{function}%
\begin{itemdecl}
~function();
\end{itemdecl}

\begin{itemdescr}
\pnum
\effects
If \tcode{*this != nullptr}, destroys the target of \keyword{this}.
\end{itemdescr}

\rSec4[func.wrap.func.mod]{Modifiers}

\indexlibrarymember{swap}{function}%
\begin{itemdecl}
void swap(function& other) noexcept;
\end{itemdecl}

\begin{itemdescr}
\pnum
\effects
Interchanges the target objects of \tcode{*this} and \tcode{other}.
\end{itemdescr}

\rSec4[func.wrap.func.cap]{Capacity}

\indexlibrarymember{operator bool}{function}%
\begin{itemdecl}
explicit operator bool() const noexcept;
\end{itemdecl}

\begin{itemdescr}
\pnum
\returns
\tcode{true} if \tcode{*this} has a target, otherwise \tcode{false}.
\end{itemdescr}

\rSec4[func.wrap.func.inv]{Invocation}

\indexlibrary{\idxcode{function}!invocation}%
\indexlibrarymember{operator()}{function}%
\begin{itemdecl}
R operator()(ArgTypes... args) const;
\end{itemdecl}

\begin{itemdescr}
\pnum
\returns
\tcode{\placeholdernc{INVOKE}<R>(f, std::forward<ArgTypes>(args)...)}\iref{func.require},
where \tcode{f} is the target object\iref{func.def} of \tcode{*this}.

\pnum
\throws
\tcode{bad_function_call} if \tcode{!*this}; otherwise, any
exception thrown by the target object.
\end{itemdescr}

\rSec4[func.wrap.func.targ]{Target access}

\indexlibrarymember{target_type}{function}%
\begin{itemdecl}
const type_info& target_type() const noexcept;
\end{itemdecl}

\begin{itemdescr}
\pnum
\returns
If \tcode{*this} has a target of type \tcode{T},
  \tcode{typeid(T)}; otherwise, \tcode{typeid(void)}.
\end{itemdescr}

\indexlibrarymember{target}{function}%
\begin{itemdecl}
template<class T>       T* target() noexcept;
template<class T> const T* target() const noexcept;
\end{itemdecl}

\begin{itemdescr}
\pnum
\returns
If \tcode{target_type() == typeid(T)}
a pointer to the stored function target; otherwise a null pointer.
\end{itemdescr}

\rSec4[func.wrap.func.nullptr]{Null pointer comparison operator functions}

\indexlibrarymember{operator==}{function}%
\begin{itemdecl}
template<class R, class... ArgTypes>
  bool operator==(const function<R(ArgTypes...)>& f, nullptr_t) noexcept;
\end{itemdecl}

\begin{itemdescr}
\pnum
\returns
\tcode{!f}.
\end{itemdescr}

\rSec4[func.wrap.func.alg]{Specialized algorithms}

\indexlibrarymember{swap}{function}%
\begin{itemdecl}
template<class R, class... ArgTypes>
  void swap(function<R(ArgTypes...)>& f1, function<R(ArgTypes...)>& f2) noexcept;
\end{itemdecl}

\begin{itemdescr}
\pnum
\effects
As if by: \tcode{f1.swap(f2);}
\end{itemdescr}%

\rSec3[func.wrap.move]{Move-only wrapper}

\rSec4[func.wrap.move.general]{General}

\pnum
The header provides partial specializations of \tcode{move_only_function}
for each combination of the possible replacements
of the placeholders \cv{}, \placeholder{ref}, and \placeholder{noex} where
\begin{itemize}
\item
\cv{} is either const or empty,
\item
\placeholder{ref} is either \tcode{\&}, \tcode{\&\&}, or empty, and
\item
\placeholder{noex} is either \tcode{true} or \tcode{false}.
\end{itemize}

\pnum
For each of the possible combinations of the placeholders mentioned above,
there is a placeholder \placeholder{inv-quals} defined as follows:
\begin{itemize}
\item
If \placeholder{ref} is empty, let \placeholder{inv-quals} be \cv{}\tcode{\&},
\item
otherwise, let \placeholder{inv-quals} be \cv{} \placeholder{ref}.
\end{itemize}

\rSec4[func.wrap.move.class]{Class template \tcode{move_only_function}}

\indexlibraryglobal{move_only_function}%
\begin{codeblock}
namespace std {
  template<class R, class... ArgTypes>
  class move_only_function<R(ArgTypes...) @\cv{}@ @\placeholder{ref}@ noexcept(@\placeholder{noex}@)> {
  public:
    using result_type = R;

    // \ref{func.wrap.move.ctor}, constructors, assignment, and destructor
    move_only_function() noexcept;
    move_only_function(nullptr_t) noexcept;
    move_only_function(move_only_function&&) noexcept;
    template<class F> move_only_function(F&&);
    template<class T, class... Args>
      explicit move_only_function(in_place_type_t<T>, Args&&...);
    template<class T, class U, class... Args>
      explicit move_only_function(in_place_type_t<T>, initializer_list<U>, Args&&...);

    move_only_function& operator=(move_only_function&&);
    move_only_function& operator=(nullptr_t) noexcept;
    template<class F> move_only_function& operator=(F&&);

    ~move_only_function();

    // \ref{func.wrap.move.inv}, invocation
    explicit operator bool() const noexcept;
    R operator()(ArgTypes...) @\cv{}@ @\placeholder{ref}@ noexcept(@\placeholder{noex}@);

    // \ref{func.wrap.move.util}, utility
    void swap(move_only_function&) noexcept;
    friend void swap(move_only_function&, move_only_function&) noexcept;
    friend bool operator==(const move_only_function&, nullptr_t) noexcept;

  private:
    template<class VT>
      static constexpr bool @\exposid{is-callable-from}@ = @\seebelow@;       // \expos
  };
}
\end{codeblock}

\pnum
The \tcode{move_only_function} class template provides polymorphic wrappers
that generalize the notion of a callable object\iref{func.def}.
These wrappers can store, move, and call arbitrary callable objects,
given a call signature.

\pnum
\recommended
Implementations should avoid the use of dynamically allocated memory
for a small contained value.
\begin{note}
Such small-object optimization can only be applied to a type \tcode{T}
for which \tcode{is_nothrow_move_constructible_v<T>} is \tcode{true}.
\end{note}

\rSec4[func.wrap.move.ctor]{Constructors, assignment, and destructor}

\indextext{move_only_function::is-callable-from@\tcode{move_only_function::\exposid{is-callable-from}}}%
\begin{itemdecl}
template<class VT>
  static constexpr bool @\exposid{is-callable-from}@ = @\seebelow@;
\end{itemdecl}

\begin{itemdescr}
\pnum
If \placeholder{noex} is \tcode{true},
\tcode{\exposid{is-callable-from}<VT>} is equal to:
\begin{codeblock}
is_nothrow_invocable_r_v<R, VT @\cv{}@ @\placeholder{ref}@, ArgTypes...> &&
is_nothrow_invocable_r_v<R, VT @\placeholder{inv-quals}@, ArgTypes...>
\end{codeblock}
Otherwise, \tcode{\exposid{is-callable-from}<VT>} is equal to:
\begin{codeblock}
is_invocable_r_v<R, VT @\cv{}@ @\placeholder{ref}@, ArgTypes...> &&
is_invocable_r_v<R, VT @\placeholder{inv-quals}@, ArgTypes...>
\end{codeblock}
\end{itemdescr}

\indexlibraryctor{move_only_function}%
\begin{itemdecl}
move_only_function() noexcept;
move_only_function(nullptr_t) noexcept;
\end{itemdecl}

\begin{itemdescr}
\pnum
\ensures
\tcode{*this} has no target object.
\end{itemdescr}

\indexlibraryctor{move_only_function}%
\begin{itemdecl}
move_only_function(move_only_function&& f) noexcept;
\end{itemdecl}

\begin{itemdescr}
\pnum
\ensures
The target object of \tcode{*this} is
the target object \tcode{f} had before construction, and
\tcode{f} is in a valid state with an unspecified value.
\end{itemdescr}

\indexlibraryctor{move_only_function}%
\begin{itemdecl}
template<class F> move_only_function(F&& f);
\end{itemdecl}

\begin{itemdescr}
\pnum
Let \tcode{VT} be \tcode{decay_t<F>}.

\pnum
\constraints
\begin{itemize}
\item
\tcode{remove_cvref_t<F>} is not the same type as \tcode{move_only_function}, and
\item
\tcode{remove_cvref_t<F>} is not a specialization of \tcode{in_place_type_t}, and
\item
\tcode{\exposid{is-callable-from}<VT>} is \tcode{true}.
\end{itemize}

\pnum
\mandates
\tcode{is_constructible_v<VT, F>} is \tcode{true}.

\pnum
\expects
\tcode{VT} meets the \oldconcept{Destructible} requirements, and
if \tcode{is_move_constructible_v<VT>} is \tcode{true},
\tcode{VT} meets the \oldconcept{MoveConstructible} requirements.

\pnum
\ensures
\tcode{*this} has no target object if any of the following hold:
\begin{itemize}
\item
\tcode{f} is a null function pointer value, or
\item
\tcode{f} is a null member pointer value, or
\item
\tcode{remove_cvref_t<F>} is a specialization of
the \tcode{move_only_function} class template,
and \tcode{f} has no target object.
\end{itemize}
Otherwise, \tcode{*this} has a target object of type \tcode{VT}
direct-non-list-initialized with \tcode{std::forward<F>(f)}.

\pnum
\throws
Any exception thrown by the initialization of the target object.
May throw \tcode{bad_alloc} unless \tcode{VT} is
a function pointer or a specialization of \tcode{reference_wrapper}.
\end{itemdescr}

\indexlibraryctor{move_only_function}%
\begin{itemdecl}
template<class T, class... Args>
  explicit move_only_function(in_place_type_t<T>, Args&&... args);
\end{itemdecl}

\begin{itemdescr}
\pnum
Let \tcode{VT} be \tcode{decay_t<T>}.

\pnum
\constraints
\begin{itemize}
\item
\tcode{is_constructible_v<VT, Args...>} is \tcode{true}, and
\item
\tcode{\exposid{is-callable-from}<VT>} is \tcode{true}.
\end{itemize}

\pnum
\mandates
\tcode{VT} is the same type as \tcode{T}.

\pnum
\expects
\tcode{VT} meets the \oldconcept{Destructible} requirements, and
if \tcode{is_move_constructible_v<VT>} is \tcode{true},
\tcode{VT} meets the \oldconcept{MoveConstructible} requirements.

\pnum
\ensures
\tcode{*this} has a target object of type \tcode{VT}
direct-non-list-initialized with \tcode{std::forward<Args>\brk{}(args)...}.

\pnum
\throws
Any exception thrown by the initialization of the target object.
May throw \tcode{bad_alloc} unless \tcode{VT} is
a function pointer or a specialization of \tcode{reference_wrapper}.
\end{itemdescr}

\indexlibraryctor{move_only_function}%
\begin{itemdecl}
template<class T, class U, class... Args>
  explicit move_only_function(in_place_type_t<T>, initializer_list<U> ilist, Args&&... args);
\end{itemdecl}

\begin{itemdescr}
\pnum
Let \tcode{VT} be \tcode{decay_t<T>}.

\pnum
\constraints
\begin{itemize}
\item
\tcode{is_constructible_v<VT, initializer_list<U>\&, Args...>} is
\tcode{true}, and
\item
\tcode{\exposid{is-callable-from}<VT>} is \tcode{true}.
\end{itemize}

\pnum
\mandates
\tcode{VT} is the same type as \tcode{T}.

\pnum
\expects
\tcode{VT} meets the \oldconcept{Destructible} requirements, and
if \tcode{is_move_constructible_v<VT>} is \tcode{true},
\tcode{VT} meets the \oldconcept{MoveConstructible} requirements.

\pnum
\ensures
\tcode{*this} has a target object of type \tcode{VT}
direct-non-list-initialized with
\tcode{ilist, std::for\-ward<Args>(args)...}.

\pnum
\throws
Any exception thrown by the initialization of the target object.
May throw \tcode{bad_alloc} unless \tcode{VT} is
a function pointer or a specialization of \tcode{reference_wrapper}.
\end{itemdescr}

\indexlibrarymember{operator=}{move_only_function}%
\begin{itemdecl}
move_only_function& operator=(move_only_function&& f);
\end{itemdecl}

\begin{itemdescr}
\pnum
\effects
Equivalent to: \tcode{move_only_function(std::move(f)).swap(*this);}

\pnum
\returns
\tcode{*this}.
\end{itemdescr}

\indexlibrarymember{operator=}{move_only_function}%
\begin{itemdecl}
move_only_function& operator=(nullptr_t) noexcept;
\end{itemdecl}

\begin{itemdescr}
\pnum
\effects
Destroys the target object of \tcode{*this}, if any.

\pnum
\returns
\tcode{*this}.
\end{itemdescr}

\indexlibrarymember{operator=}{move_only_function}%
\begin{itemdecl}
template<class F> move_only_function& operator=(F&& f);
\end{itemdecl}

\begin{itemdescr}
\pnum
\effects
Equivalent to: \tcode{move_only_function(std::forward<F>(f)).swap(*this);}

\pnum
\returns
\tcode{*this}.
\end{itemdescr}

\indexlibrarydtor{move_only_function}%
\begin{itemdecl}
~move_only_function();
\end{itemdecl}

\begin{itemdescr}
\pnum
\effects
Destroys the target object of \tcode{*this}, if any.
\end{itemdescr}

\rSec4[func.wrap.move.inv]{Invocation}

\indexlibrarymember{operator bool}{move_only_function}%
\begin{itemdecl}
explicit operator bool() const noexcept;
\end{itemdecl}

\begin{itemdescr}
\pnum
\returns
\tcode{true} if \tcode{*this} has a target object, otherwise \tcode{false}.
\end{itemdescr}

\indexlibrarymember{operator()}{move_only_function}%
\begin{itemdecl}
R operator()(ArgTypes... args) @\cv{}@ @\placeholder{ref}@ noexcept(@\placeholder{noex}@);
\end{itemdecl}

\begin{itemdescr}
\pnum
\expects
\tcode{*this} has a target object.

\pnum
\effects
Equivalent to:
\begin{codeblock}
return @\placeholder{INVOKE}@<R>(static_cast<F @\placeholder{inv-quals}@>(f), std::forward<ArgTypes>(args)...);
\end{codeblock}
where \tcode{f} is an lvalue designating the target object of \tcode{*this} and
\tcode{F} is the type of \tcode{f}.
\end{itemdescr}

\rSec4[func.wrap.move.util]{Utility}

\indexlibrarymember{swap}{move_only_function}%
\begin{itemdecl}
void swap(move_only_function& other) noexcept;
\end{itemdecl}

\begin{itemdescr}
\pnum
\effects
Exchanges the target objects of \tcode{*this} and \tcode{other}.
\end{itemdescr}

\indexlibrarymember{swap}{move_only_function}%
\begin{itemdecl}
friend void swap(move_only_function& f1, move_only_function& f2) noexcept;
\end{itemdecl}

\begin{itemdescr}
\pnum
\effects
Equivalent to \tcode{f1.swap(f2)}.
\end{itemdescr}

\indexlibrarymember{operator==}{move_only_function}%
\begin{itemdecl}
friend bool operator==(const move_only_function& f, nullptr_t) noexcept;
\end{itemdecl}

\begin{itemdescr}
\pnum
\returns
\tcode{true} if \tcode{f} has no target object, otherwise \tcode{false}.
\end{itemdescr}

\rSec3[func.wrap.copy]{Copyable wrapper}

\rSec4[func.wrap.copy.general]{General}

\pnum
The header provides partial specializations of \tcode{copyable_function}
for each combination of the possible replacements
of the placeholders \cv{}, \placeholder{ref}, and \placeholder{noex} where
\begin{itemize}
\item
\cv{} is either const or empty,
\item
\placeholder{ref} is either \tcode{\&}, \tcode{\&\&}, or empty, and
\item
\placeholder{noex} is either \tcode{true} or \tcode{false}.
\end{itemize}

\pnum
For each of the possible combinations of the placeholders mentioned above,
there is a placeholder \placeholder{inv-quals} defined as follows:
\begin{itemize}
\item
If \placeholder{ref} is empty, let \placeholder{inv-quals} be \cv{}\tcode{\&},
\item
otherwise, let \placeholder{inv-quals} be \cv{} \placeholder{ref}.
\end{itemize}

\rSec4[func.wrap.copy.class]{Class template \tcode{copyable_function}}

\indexlibraryglobal{copyable_function}%
\begin{codeblock}
namespace std {
  template<class R, class... ArgTypes>
  class copyable_function<R(ArgTypes...) @\cv{}@ @\placeholder{ref}@ noexcept(@\placeholder{noex}@)> {
  public:
    using result_type = R;

    // \ref{func.wrap.copy.ctor}, constructors, assignments, and destructors
    copyable_function() noexcept;
    copyable_function(nullptr_t) noexcept;
    copyable_function(const copyable_function&);
    copyable_function(copyable_function&&) noexcept;
    template<class F> copyable_function(F&&);
    template<class T, class... Args>
      explicit copyable_function(in_place_type_t<T>, Args&&...);
    template<class T, class U, class... Args>
      explicit copyable_function(in_place_type_t<T>, initializer_list<U>, Args&&...);

    copyable_function& operator=(const copyable_function&);
    copyable_function& operator=(copyable_function&&);
    copyable_function& operator=(nullptr_t) noexcept;
    template<class F> copyable_function& operator=(F&&);

    ~copyable_function();

    // \ref{func.wrap.copy.inv}, invocation
    explicit operator bool() const noexcept;
    R operator()(ArgTypes...) @\cv{}@ @\placeholder{ref}@ noexcept(@\placeholder{noex}@);

    // \ref{func.wrap.copy.util}, utility
    void swap(copyable_function&) noexcept;
    friend void swap(copyable_function&, copyable_function&) noexcept;
    friend bool operator==(const copyable_function&, nullptr_t) noexcept;

  private:
    template<class VT>
      static constexpr bool @\exposid{is-callable-from}@ = @\seebelow@;       // \expos
  };
}
\end{codeblock}

\pnum
The \tcode{copyable_function} class template provides polymorphic wrappers
that generalize the notion of a callable object\iref{func.def}.
These wrappers can store, copy, move, and call arbitrary callable objects,
given a call signature.

\pnum
\recommended
Implementations should avoid the use of dynamically allocated memory
for a small contained value.
\begin{note}
Such small-object optimization can only be applied to a type \tcode{T}
for which \tcode{is_nothrow_move_constructible_v<T>} is \tcode{true}.
\end{note}

\rSec4[func.wrap.copy.ctor]{Constructors, assignments, and destructors}

\indextext{copyable_function::is-callable-from@\tcode{copyable_function::\exposid{is-callable-from}}}%
\begin{itemdecl}
template<class VT>
  static constexpr bool @\exposid{is-callable-from}@ = @\seebelow@;
\end{itemdecl}

\begin{itemdescr}
\pnum
If \placeholder{noex} is \tcode{true},
\tcode{\exposid{is-callable-from}<VT>} is equal to:
\begin{codeblock}
is_nothrow_invocable_r_v<R, VT @\cv{}@ @\placeholder{ref}@, ArgTypes...> &&
is_nothrow_invocable_r_v<R, VT @\placeholder{inv-quals}@, ArgTypes...>
\end{codeblock}
Otherwise, \tcode{\exposid{is-callable-from}<VT>} is equal to:
\begin{codeblock}
is_invocable_r_v<R, VT @\cv{}@ @\placeholder{ref}@, ArgTypes...> &&
is_invocable_r_v<R, VT @\placeholder{inv-quals}@, ArgTypes...>
\end{codeblock}
\end{itemdescr}

\indexlibraryctor{copyable_function}%
\begin{itemdecl}
copyable_function() noexcept;
copyable_function(nullptr_t) noexcept;
\end{itemdecl}

\begin{itemdescr}
\pnum
\ensures
\tcode{*this} has no target object.
\end{itemdescr}

\indexlibraryctor{copyable_function}%
\begin{itemdecl}
copyable_function(const copyable_function& f);
\end{itemdecl}

\begin{itemdescr}
\pnum
\ensures
\tcode{*this} has no target object if \tcode{f} had no target object.
Otherwise, the target object of \tcode{*this}
is a copy of the target object of \tcode{f}.

\pnum
\throws
Any exception thrown by the initialization of the target object.
May throw \tcode{bad_alloc}.
\end{itemdescr}

\indexlibraryctor{copyable_function}%
\begin{itemdecl}
copyable_function(copyable_function&& f) noexcept;
\end{itemdecl}

\begin{itemdescr}
\pnum
\ensures
The target object of \tcode{*this} is
the target object \tcode{f} had before construction, and
\tcode{f} is in a valid state with an unspecified value.
\end{itemdescr}

\indexlibraryctor{copyable_function}%
\begin{itemdecl}
template<class F> copyable_function(F&& f);
\end{itemdecl}

\begin{itemdescr}
\pnum
Let \tcode{VT} be \tcode{decay_t<F>}.

\pnum
\constraints
\begin{itemize}
\item
\tcode{remove_cvref_t<F>} is not the same type as \tcode{copyable_function}, and
\item
\tcode{remove_cvref_t<F>} is not a specialization of \tcode{in_place_type_t}, and
\item
\tcode{\exposid{is-callable-from}<VT>} is \tcode{true}.
\end{itemize}

\pnum
\mandates
\begin{itemize}
\item
\tcode{is_constructible_v<VT, F>} is \tcode{true}, and
\item
\tcode{is_copy_constructible_v<VT>} is \tcode{true}.
\end{itemize}

\pnum
\expects
\tcode{VT} meets the \oldconcept{Destructible} and
\oldconcept{CopyConstructible} requirements.

\pnum
\ensures
\tcode{*this} has no target object if any of the following hold:
\begin{itemize}
\item
\tcode{f} is a null function pointer value, or
\item
\tcode{f} is a null member pointer value, or
\item
\tcode{remove_cvref_t<F>} is a specialization of
the \tcode{copyable_function} class template,
and \tcode{f} has no target object.
\end{itemize}
Otherwise, \tcode{*this} has a target object of type \tcode{VT}
direct-non-list-initialized with \tcode{std::forward<F>(f)}.

\pnum
\throws
Any exception thrown by the initialization of the target object.
May throw \tcode{bad_alloc} unless \tcode{VT} is
a function pointer or a specialization of \tcode{reference_wrapper}.
\end{itemdescr}

\indexlibraryctor{copyable_function}%
\begin{itemdecl}
template<class T, class... Args>
  explicit copyable_function(in_place_type_t<T>, Args&&... args);
\end{itemdecl}

\begin{itemdescr}
\pnum
Let \tcode{VT} be \tcode{decay_t<T>}.

\pnum
\constraints
\begin{itemize}
\item
\tcode{is_constructible_v<VT, Args...>} is \tcode{true}, and
\item
\tcode{\exposid{is-callable-from}<VT>} is \tcode{true}.
\end{itemize}

\pnum
\mandates
\begin{itemize}
\item
\tcode{VT} is the same type as \tcode{T}, and
\item
\tcode{is_copy_constructible_v<VT>} is \tcode{true}.
\end{itemize}

\pnum
\expects
\tcode{VT} meets the \oldconcept{Destructible} and
\oldconcept{CopyConstructible} requirements.

\pnum
\ensures
\tcode{*this} has a target object of type \tcode{VT}
direct-non-list-initialized with \tcode{std::forward<Args>\brk{}(args)...}.

\pnum
\throws
Any exception thrown by the initialization of the target object.
May throw \tcode{bad_alloc} unless \tcode{VT} is
a pointer or a specialization of \tcode{reference_wrapper}.
\end{itemdescr}

\indexlibraryctor{copyable_function}%
\begin{itemdecl}
template<class T, class U, class... Args>
  explicit copyable_function(in_place_type_t<T>, initializer_list<U> ilist, Args&&... args);
\end{itemdecl}

\begin{itemdescr}
\pnum
Let \tcode{VT} be \tcode{decay_t<T>}.

\pnum
\constraints
\begin{itemize}
\item
\tcode{is_constructible_v<VT, initializer_list<U>\&, Args...>} is
\tcode{true}, and
\item
\tcode{\exposid{is-callable-from}<VT>} is \tcode{true}.
\end{itemize}

\pnum
\mandates
\begin{itemize}
\item
\tcode{VT} is the same type as \tcode{T}, and
\item
\tcode{is_copy_constructible_v<VT>} is \tcode{true}.
\end{itemize}

\pnum
\expects
\tcode{VT} meets the \oldconcept{Destructible} and
\oldconcept{CopyConstructible} requirements.

\pnum
\ensures
\tcode{*this} has a target object of type \tcode{VT}
direct-non-list-initialized with
\tcode{ilist, std::for\-ward<Args>(args)...}.

\pnum
\throws
Any exception thrown by the initialization of the target object.
May throw \tcode{bad_alloc} unless \tcode{VT} is
a pointer or a specialization of \tcode{reference_wrapper}.
\end{itemdescr}

\indexlibrarymember{operator=}{copyable_function}%
\begin{itemdecl}
copyable_function& operator=(const copyable_function& f);
\end{itemdecl}

\begin{itemdescr}
\pnum
\effects
Equivalent to: \tcode{copyable_function(f).swap(*this);}

\pnum
\returns
\tcode{*this}.
\end{itemdescr}

\indexlibrarymember{operator=}{copyable_function}%
\begin{itemdecl}
copyable_function& operator=(copyable_function&& f);
\end{itemdecl}

\begin{itemdescr}
\pnum
\effects
Equivalent to: \tcode{copyable_function(std::move(f)).swap(*this);}

\pnum
\returns
\tcode{*this}.
\end{itemdescr}

\indexlibrarymember{operator=}{copyable_function}%
\begin{itemdecl}
copyable_function& operator=(nullptr_t) noexcept;
\end{itemdecl}

\begin{itemdescr}
\pnum
\effects
Destroys the target object of \tcode{*this}, if any.

\pnum
\returns
\tcode{*this}.
\end{itemdescr}

\indexlibrarymember{operator=}{copyable_function}%
\begin{itemdecl}
template<class F> copyable_function& operator=(F&& f);
\end{itemdecl}

\begin{itemdescr}
\pnum
\effects
Equivalent to: \tcode{copyable_function(std::forward<F>(f)).swap(*this);}

\pnum
\returns
\tcode{*this}.
\end{itemdescr}

\indexlibrarydtor{copyable_function}%
\begin{itemdecl}
~copyable_function();
\end{itemdecl}

\begin{itemdescr}
\pnum
\effects
Destroys the target object of \tcode{*this}, if any.
\end{itemdescr}

\rSec4[func.wrap.copy.inv]{Invocation}

\indexlibrarymember{operator bool}{copyable_function}%
\begin{itemdecl}
explicit operator bool() const noexcept;
\end{itemdecl}

\begin{itemdescr}
\pnum
\returns
\tcode{true} if \tcode{*this} has a target object, otherwise \tcode{false}.
\end{itemdescr}

\indexlibrarymember{operator()}{copyable_function}%
\begin{itemdecl}
R operator()(ArgTypes... args) @\cv{}@ @\placeholder{ref}@ noexcept(@\placeholder{noex}@);
\end{itemdecl}

\begin{itemdescr}
\pnum
\expects
\tcode{*this} has a target object.

\pnum
\effects
Equivalent to:
\begin{codeblock}
return @\placeholder{INVOKE}@<R>(static_cast<F @\placeholder{inv-quals}@>(f), std::forward<ArgTypes>(args)...);
\end{codeblock}
where \tcode{f} is an lvalue designating the target object of \tcode{*this} and
\tcode{F} is the type of \tcode{f}.
\end{itemdescr}

\rSec4[func.wrap.copy.util]{Utility}

\indexlibrarymember{swap}{copyable_function}%
\begin{itemdecl}
void swap(copyable_function& other) noexcept;
\end{itemdecl}

\begin{itemdescr}
\pnum
\effects
Exchanges the target objects of \tcode{*this} and \tcode{other}.
\end{itemdescr}

\indexlibrarymember{swap}{copyable_function}%
\begin{itemdecl}
friend void swap(copyable_function& f1, copyable_function& f2) noexcept;
\end{itemdecl}

\begin{itemdescr}
\pnum
\effects
Equivalent to \tcode{f1.swap(f2)}.
\end{itemdescr}

\indexlibrarymember{operator==}{copyable_function}%
\begin{itemdecl}
friend bool operator==(const copyable_function& f, nullptr_t) noexcept;
\end{itemdecl}

\begin{itemdescr}
\pnum
\returns
\tcode{true} if \tcode{f} has no target object, otherwise \tcode{false}.
\end{itemdescr}

\rSec3[func.wrap.ref]{Non-owning wrapper}

\rSec4[func.wrap.ref.general]{General}

\pnum
The header provides partial specializations of \tcode{function_ref}
for each combination of the possible replacements of
the placeholders \cv{} and \placeholder{noex} where:

\begin{itemize}
\item \cv{} is either const or empty, and
\item \placeholder{noex} is either \tcode{true} or \tcode{false}.
\end{itemize}

\rSec4[func.wrap.ref.class]{Class template \tcode{function_ref}}

\indexlibraryglobal{function_ref}%
\begin{codeblock}
namespace std {
  template<class R, class... ArgTypes>
  class function_ref<R(ArgTypes...) @\cv{}@ noexcept(@\placeholder{noex}@)> {
  public:
    // \ref{func.wrap.ref.ctor}, constructors and assignment operators
    template<class F> function_ref(F*) noexcept;
    template<class F> constexpr function_ref(F&&) noexcept;
    template<auto f> constexpr function_ref(nontype_t<f>) noexcept;
    template<auto f, class U> constexpr function_ref(nontype_t<f>, U&&) noexcept;
    template<auto f, class T> constexpr function_ref(nontype_t<f>, @\cv{}@ T*) noexcept;

    constexpr function_ref(const function_ref&) noexcept = default;
    constexpr function_ref& operator=(const function_ref&) noexcept = default;
    template<class T> function_ref& operator=(T) = delete;

    // \ref{func.wrap.ref.inv}, invocation
    R operator()(ArgTypes...) const noexcept(@\placeholder{noex}@);

  private:
    template<class... T>
      static constexpr bool @\exposidnc{is-invocable-using}@ = @\seebelownc@;     // \expos

    R (*@\exposidnc{thunk-ptr}@)(@\exposidnc{BoundEntityType}@, Args&&...) noexcept(@\placeholdernc{noex}@);  // \expos
    @\exposidnc{BoundEntityType}@ @\exposidnc{bound-entity}@;                               // \expos
  };

  // \ref{func.wrap.ref.deduct}, deduction guides
  template<class F>
    function_ref(F*) -> function_ref<F>;
  template<auto f>
    function_ref(nontype_t<f>) -> function_ref<@\seebelow@>;
  template<auto f, class T>
    function_ref(nontype_t<f>, T&&) -> function_ref<@\seebelow@>;
}
\end{codeblock}

\pnum
An object of class
\tcode{function_ref<R(Args...) \cv{} noexcept(\placeholder{noex})>}
stores a pointer to function \exposid{thunk-ptr} and
an object \exposid{bound-entity}.
\exposid{bound-entity} has
an unspecified trivially copyable type \exposid{BoundEntityType}, that
models \libconcept{copyable} and
is capable of storing a pointer to object value or a pointer to function value.
The type of \exposid{thunk-ptr} is
\tcode{R(*)(\exposidnc{BoundEntityType}, Args\&\&...) noexcept(\placeholder{noex})}.

\pnum
Each specialization of \tcode{function_ref} is
a trivially copyable type\iref{term.trivially.copyable.type}
that models \libconcept{copyable}.

\pnum
Within \ref{func.wrap.ref},
\tcode{\placeholder{call-args}} is an argument pack with elements such that
\tcode{decltype((\placeholder{call-args}\linebreak{}))...} denote
\tcode{Args\&\&...} respectively.

\rSec4[func.wrap.ref.ctor]{Constructors and assignment operators}

\indextext{function_ref::is-invocable-using@\tcode{function_ref::\exposid{is-invocable-using}}}%
\begin{itemdecl}
template<class... T>
  static constexpr bool @\exposid{is-invocable-using}@ = @\seebelow@;
\end{itemdecl}

\begin{itemdescr}
\pnum
If \placeholder{noex} is \tcode{true},
\tcode{\exposid{is-invocable-using}<T...>} is equal to:
\begin{codeblock}
is_nothrow_invocable_r_v<R, T..., ArgTypes...>
\end{codeblock}
Otherwise, \tcode{\exposid{is-invocable-using}<T...>} is equal to:
\begin{codeblock}
is_invocable_r_v<R, T..., ArgTypes...>
\end{codeblock}
\end{itemdescr}

\indexlibraryctor{function_ref}%
\begin{itemdecl}
template<class F> function_ref(F* f) noexcept;
\end{itemdecl}

\begin{itemdescr}
\pnum
\constraints
\begin{itemize}
\item \tcode{is_function_v<F>} is \tcode{true}, and
\item \tcode{\exposid{is-invocable-using}<F>} is \tcode{true}.
\end{itemize}

\pnum
\expects
\tcode{f} is not a null pointer.

\pnum
\effects
Initializes
\exposid{bound-entity} with \tcode{f}, and
\exposid{thunk-ptr} with the address of a function \tcode{\placeholder{thunk}}
such that
\tcode{\placeholder{thunk}(\exposid{bound-entity}, \placeholder{call-args}...)}
is expression-equivalent\iref{defns.expression.equivalent} to
\tcode{invoke_r<R>(f, \placeholder{call-args}...)}.
\end{itemdescr}

\indexlibraryctor{function_ref}%
\begin{itemdecl}
template<class F> constexpr function_ref(F&& f) noexcept;
\end{itemdecl}

\begin{itemdescr}
\pnum
Let \tcode{T} be \tcode{remove_reference_t<F>}.

\pnum
\constraints
\begin{itemize}
\item \tcode{remove_cvref_t<F>} is not the same type as \tcode{function_ref},
\item \tcode{is_member_pointer_v<T>} is \tcode{false}, and
\item \tcode{\exposid{is-invocable-using}<\cv{} T\&>} is \tcode{true}.
\end{itemize}

\pnum
\effects
Initializes
\exposid{bound-entity} with \tcode{addressof(f)}, and
\exposid{thunk-ptr} with the address of a function \tcode{\placeholder{thunk}}
such that
\tcode{\placeholder{thunk}(\exposid{bound-entity}, \placeholder{call-args}...)}
is expression-equivalent\iref{defns.expression.equivalent} to
\tcode{invoke_r<R>(static_cast<\cv{} T\&>(f), \placeholder{call-args}...)}.
\end{itemdescr}

\indexlibraryctor{function_ref}%
\begin{itemdecl}
template<auto f> constexpr function_ref(nontype_t<f>) noexcept;
\end{itemdecl}

\begin{itemdescr}
\pnum
Let \tcode{F} be \tcode{decltype(f)}.

\pnum
\constraints
\tcode{\exposid{is-invocable-using}<F>} is \tcode{true}.

\pnum
\mandates
If \tcode{is_pointer_v<F> || is_member_pointer_v<F>} is \tcode{true},
then \tcode{f != nullptr} is \tcode{true}.

\pnum
\effects
Initializes
\exposid{bound-entity} with a pointer to an unspecified object or
null pointer value, and
\exposid{thunk-ptr} with the address of a function \tcode{\placeholder{thunk}}
such that
\tcode{\placeholder{thunk}(\exposid{bound-entity}, \placeholder{call-args}...)}
is expression-equivalent\iref{defns.expression.equivalent} to
\tcode{invoke_r<R>(f, \placeholder{call-args}...)}.
\end{itemdescr}

\indexlibraryctor{function_ref}%
\begin{itemdecl}
template<auto f, class U>
  constexpr function_ref(nontype_t<f>, U&& obj) noexcept;
\end{itemdecl}

\begin{itemdescr}
\pnum
Let \tcode{T} be \tcode{remove_reference_t<U>} and
\tcode{F} be \tcode{decltype(f)}.

\pnum
\constraints
\begin{itemize}
\item \tcode{is_rvalue_reference_v<U\&\&>} is \tcode{false}, and
\item \tcode{\exposid{is-invocable-using}<F, \cv{} T\&>} is \tcode{true}.
\end{itemize}

\pnum
\mandates
If \tcode{is_pointer_v<F> || is_member_pointer_v<F>} is \tcode{true},
then \tcode{f != nullptr} is \tcode{true}.

\pnum
\effects
Initializes
\exposid{bound-entity} with \tcode{addressof(obj)}, and
\exposid{thunk-ptr} with the address of a function \tcode{\placeholder{thunk}}
such that
\tcode{\placeholder{thunk}(\exposid{bound-entity}, \placeholder{call-args}...)}
is expression-equivalent\iref{defns.expression.equivalent} to
\tcode{invoke_r<R>(f, static_cast<\cv{} T\&>(obj), \placeholder{call-args}...)}.
\end{itemdescr}

\indexlibraryctor{function_ref}%
\begin{itemdecl}
template<auto f, class T>
  constexpr function_ref(nontype_t<f>, @\cv{}@ T* obj) noexcept;
\end{itemdecl}

\begin{itemdescr}
\pnum
Let \tcode{F} be \tcode{decltype(f)}.

\pnum
\constraints
\tcode{\exposid{is-invocable-using}<F, \cv{} T*>} is \tcode{true}.

\pnum
\mandates
If \tcode{is_pointer_v<F> || is_member_pointer_v<F>} is \tcode{true},
then \tcode{f != nullptr} is \tcode{true}.

\pnum
\expects
If \tcode{is_member_pointer_v<F>} is \tcode{true},
\tcode{obj} is not a null pointer.

\pnum
\effects
Initializes
\exposid{bound-entity} with \tcode{obj}, and
\exposid{thunk-ptr} with the address of a function \tcode{\placeholder{thunk}}
such that
\tcode{\placeholder{thunk}(\exposid{bound-entity}, \placeholder{call-args}...)}
is expression-equivalent\iref{defns.expression.equivalent} to
\tcode{invoke_r<R>(f, obj, \placeholder{call-args}...)}.
\end{itemdescr}

\indexlibrarymember{operator=}{function_ref}%
\begin{itemdecl}
template<class T> function_ref& operator=(T) = delete;
\end{itemdecl}

\begin{itemdescr}
\pnum
\constraints
\begin{itemize}
\item \tcode{T} is not the same type as \tcode{function_ref},
\item \tcode{is_pointer_v<T>} is \tcode{false}, and
\item \tcode{T} is not a specialization of \tcode{nontype_t}.
\end{itemize}
\end{itemdescr}

\rSec4[func.wrap.ref.inv]{Invocation}

\indexlibrarymember{operator()}{function_ref}%
\begin{itemdecl}
R operator()(ArgTypes... args) const noexcept(@\placeholder{noex}@);
\end{itemdecl}

\begin{itemdescr}
\pnum
\effects
Equivalent to:
\tcode{return \exposid{thunk-ptr}(\exposid{bound-entity}, std::forward<ArgTypes>(args)...);}
\end{itemdescr}

\rSec4[func.wrap.ref.deduct]{Deduction guides}

\begin{itemdecl}
template<class F>
  function_ref(F*) -> function_ref<F>;
\end{itemdecl}

\begin{itemdescr}
\pnum
\constraints
\tcode{is_function_v<F>} is \tcode{true}.
\end{itemdescr}

\begin{itemdecl}
template<auto f>
  function_ref(nontype_t<f>) -> function_ref<@\seebelow@>;
\end{itemdecl}

\begin{itemdescr}
\pnum
Let \tcode{F} be \tcode{remove_pointer_t<decltype(f)>}.

\pnum
\constraints
\tcode{is_function_v<F>} is \tcode{true}.

\pnum
\remarks
The deduced type is \tcode{function_ref<F>}.
\end{itemdescr}

\begin{itemdecl}
template<auto f, class T>
  function_ref(nontype_t<f>, T&&) -> function_ref<@\seebelow@>;
\end{itemdecl}

\begin{itemdescr}
\pnum
Let \tcode{F} be \tcode{decltype(f)}.

\pnum
\constraints
%FIXME: R and E should be defined outside of these constraints.
%FIXME: Define R and E via "let" in paragraph above, then use them here and below.
\begin{itemize}
\item
\tcode{F} is of the form
\tcode{R(G::*)(A...) \cv{} \opt{\&} noexcept(E)} for a type \tcode{G}, or
\item
\tcode{F} is of the form
\tcode{M G::*} for a type \tcode{G} and an object type \tcode{M},
in which case
let \tcode{R} be \tcode{invoke_result_t<F, T\&>},
\tcode{A...} be an empty pack, and
\tcode{E} be \tcode{false}, or
\item
\tcode{F} is of the form
\tcode{R(*)(G, A...) noexcept(E)} for a type \tcode{G}.
\end{itemize}

\pnum
\remarks
The deduced type is \tcode{function_ref<R(A...) noexcept(E)>}.
\end{itemdescr}
\indextext{function object!wrapper|)}
